\section{Conclusions}
\label{sec:conclusions}

Project~2 establishes a physics-driven Falcon~9/Dragon simulator and a sensor-driven navigation system for the selected CRS-10 scenario. Forces and moments generate motion through six-degree-of-freedom equations; changing propellant inventories determine mass, CG and inertia; configuration events preserve inherited child states. The simulator therefore generates the trajectory from its models and commands rather than replaying a prescribed flight record. The shared conventions, module interfaces and reproducible execution provide the computational foundation required for subsequent guidance and control development \cite[slides 2--4]{project2}.

The demonstrated mission connects ascent, shutdown, separation, booster return and contact, while independently continuing stage~2 and Dragon to $600\,\mathrm{s}$. The model/navigation and environment suites total 171 passing assertion-based tests. The separation audit satisfies the selected conservation criteria, and the nominal booster meets all declared contact limits. The isolated quaternion fixture exhibits fourth-order RK4 convergence. These checks verify the implemented scenario and interfaces, rather than certify real Falcon~9 separation or landing performance.

The public ascent comparison quantifies the available external evidence: altitude RMSE is $1.577\,\mathrm{km}$ and speed RMSE is $18.46\,\mathrm{m\,s^{-1}}$ over the retained interval. Because this record also informs the ascent calibration, it is not an independent validation set. The return sequence, upper-stage commands and internal mass distribution remain evidence-qualified research models rather than reconstructed CRS-10 flight software or measured mission internals.

Navigation propagates noisy gyro measurements and its own attitude, aided by GNSS, valid LiDAR and 381 accepted star fixes. The nominal IMU-station errors are $\PositionRMS\,\mathrm{m}$ position RMS, $\VelocityRMS\,\mathrm{m\,s^{-1}}$ velocity RMS and $\AttitudeRMS^\circ$ attitude RMS. Final attitude error is $\AttitudeFinal^\circ$, compared with $\GyroOnlyFinal^\circ$ for the paired unaided quaternion. Independent nonzero-error fixtures demonstrate accelerometer/gyro calibration within stated tolerances, and 20 independent realizations pass the position-marginal ANEES test. Nominal calibrated instantaneous rate RMS remains slightly above raw RMS; calibration benefit is demonstrated in the separate nonzero-error fixture. Accuracy, consistency, aiding and parameter observability therefore remain distinct claims.

The principal engineering outcome is a traceable environment in which those distinctions can be investigated. Mixed-source mass properties, assumed return aerodynamics, calm nominal wind, ideal moments and truth-state command feedback limit extrapolation to integrated GNC. Physical actuator modeling, broader uncertainty/outage trials and independent flight validation are justified next steps before closing the navigation-to-control loop. Within this boundary, the project delivers the vehicle, environment, measurement, estimation and numerical interfaces, quantified tests and an auditable assumption/source register. The archived raw evidence and executable figure/validation scripts make the reported results reproducible from a fresh checkout.
